% TOP_PROBLEM: {"id": 30001856, "title": "Higher-Dimensional Riesz Transforms and Uniform Rectifiability", "category_id": 8, "category": {"id": 8, "name": "analysis", "display_name": "Analysis", "description": "Limits, continuity, calculus, and function theory.", "slug": "analysis", "order_index": 8, "created_at": "2026-07-31T15:26:25.671Z"}, "status": "open"}
% FIRST_POSTED: 2026-10-06
% ATTEMPT_STATUS: unresolved
% OPENAI_CLAIM: {"schema_version":1,"source_repo":"https://github.com/openai/math","source_commit":"adc7f1241b42e322a6451854ab7e4b4c146bf78a","release_date":"2026-10-06","match":"partial","resolution":"partial","independently_reviewed":false,"second_pass":"n/a","families":[{"family":"081","title":"Riesz transforms and rectifiability in higher codimension","subject":"Real and complex analysis","manuscripts":[{"title":"Riesz transforms and uniform rectifiability in higher codimension","pdf_path":"preprints/Riesz-transforms-and-uniform-rectifiability-in-higher-codimension-September-24-2026/paper.pdf","readme_path":"preprints/Riesz-transforms-and-uniform-rectifiability-in-higher-codimension-September-24-2026/README.md","date":"2026-09-24","theorem_ref":"Theorem/corollary (label thm:main), preprints/Riesz-transforms-and-uniform-rectifiability-in-higher-codimension-September-24-2026/build/sections/00-introduction.tex"}],"lean":{"doc_path":"lean/docs/081.md","comparators":["lean/ComparatorChallenges/RieszQuantitative.json","lean/ComparatorChallenges/RieszQuantitative.lean"],"declarations":["OAI.RieszRectifiability.quantitative_higher_codimension_riesz_rectifiability"],"covers_main_theorem":true},"reasoning_trace":null}]}
\documentclass[11pt]{article}
\usepackage[margin=25mm]{geometry}
\usepackage{fontspec}
\setmainfont{Latin Modern Roman}
\usepackage{amsmath,amssymb,mathtools}
\usepackage{unicode-math}
\setmathfont{Latin Modern Math}
\usepackage{enumitem}
\usepackage{xurl,hyperref,bookmark}
\hypersetup{colorlinks=true,urlcolor=blue}
\setcounter{secnumdepth}{0}
\setlength{\parindent}{0pt}
\setlength{\parskip}{5pt}
\setlength{\emergencystretch}{3em}
\begin{document}
\section*{Higher-Dimensional Riesz Transforms and Uniform Rectifiability}

\subsection{OpenAI's claim}
\textit{Higher-codimension Riesz rectifiability}: The claimed result turns uniform $L^2$ control of all positive hard truncations of the Riesz transform into quantitative uniform rectifiability. It concerns $n$-regular measures in $\mathbb R^d$ with $d\ge4$ and $2\le n\le d-2$.

\subsection{Outline of the argument}
\textit{Higher-codimension Riesz rectifiability}: After a flat blowup, the normal height is renormalized by the initial excess so that relative improvement can be detected. A fractional-energy test yields compactness and a limiting equation whose integrability forces the height to be affine. Dyadic packing then counts failures of excess improvement, leading to large Lipschitz pieces at every support scale.

\subsection{Scope relative to this problem}
\textit{Higher-codimension Riesz rectifiability}: The main theorem proves the Riesz-bound implication in all intermediate dimensions 2<=n<=d-2, and the consequences give variation/maximal control there. The displayed question quantifies more broadly over other integer dimensions (including n=1 and/or codimension one), which the main statement does not itself cover; no full label is inferred for this broader record.

This page records partial support. Complete resolution of the catalogue question is not established by this record.

\subsection{Formal verification}
Family 081 describes Lean coverage for the quantitative Lipschitz-piece conclusion in the stated higher-codimension Riesz setting.

\textit{Higher-codimension Riesz rectifiability}: The manuscript links \texttt{RieszQuantitative.lean} (declaration \texttt{OAI.RieszRectifiability.quantitative\_higher\_codimension\_riesz\_rectifiability}).

The documented selection includes this manuscript’s main assertion; that assertion retains the restrictions stated above.

This site has not run the Lean development. These links describe OpenAI’s selected formal statements and should be read with their stated scope.

\subsection{Status caveat}
OpenAI reports that an internal model produced the manuscripts. Their arguments have not been independently verified as proofs by this site. Any portion outside the documented formal scope may contain mistakes, and the catalogue label reports the announced mathematical claim.

\subsection{Sources}
\begin{itemize}
\item OpenAI, \emph{Riesz transforms and uniform rectifiability in higher codimension}, 2026-09-24; citation key \texttt{OAI:Riesz-transforms-and-uniform-rectifiability-in-higher-codimension-September-24-2026}. \href{https://github.com/openai/math/blob/adc7f1241b42e322a6451854ab7e4b4c146bf78a/preprints/Riesz-transforms-and-uniform-rectifiability-in-higher-codimension-September-24-2026/paper.pdf}{PDF (GitHub)}; \href{https://github.com/openai/math/raw/adc7f1241b42e322a6451854ab7e4b4c146bf78a/preprints/Riesz-transforms-and-uniform-rectifiability-in-higher-codimension-September-24-2026/paper.pdf}{PDF (direct)}; \href{https://github.com/openai/math/blob/adc7f1241b42e322a6451854ab7e4b4c146bf78a/preprints/Riesz-transforms-and-uniform-rectifiability-in-higher-codimension-September-24-2026/README.md}{Manuscript page and citation}; \href{https://github.com/openai/math/blob/main/preprints/Riesz-transforms-and-uniform-rectifiability-in-higher-codimension-September-24-2026/paper.pdf}{Latest version}.
\item Family 081: \href{https://github.com/openai/math/blob/adc7f1241b42e322a6451854ab7e4b4c146bf78a/lean/docs/081.md}{Lean scope}.
\item \href{https://github.com/openai/math/blob/adc7f1241b42e322a6451854ab7e4b4c146bf78a/lean/ComparatorChallenges/RieszQuantitative.json}{Comparator RieszQuantitative.json}.
\item \href{https://github.com/openai/math/blob/adc7f1241b42e322a6451854ab7e4b4c146bf78a/lean/ComparatorChallenges/RieszQuantitative.lean}{Comparator RieszQuantitative.lean}.
\item \href{https://github.com/openai/math/blob/adc7f1241b42e322a6451854ab7e4b4c146bf78a/CONTENTS.md}{OpenAI catalogue entry}.
\end{itemize}
\end{document}
